\section{Description du problème}
\label{sec:probleme}

Le problème de planification de maintenance ferroviaire consiste à organiser les opérations de maintenance sur un horizon temporel donné, tout en respectant de nombreuses contraintes de ressources et de disponibilité des trains. L'objectif principal est de minimiser les retards des maintenances par rapport à leurs dates butoirs, en tenant compte des priorités.

\subsection{Entités et paramètres}

Les \textbf{trains} constituent l'élément central : chaque train possède une ou plusieurs fenêtres de présence au centre de maintenance, définies par des temps d'arrivée et de départ. Les trains sont caractérisés par leur type, qui détermine leurs propriétés techniques (étage requis, complexité influençant les temps de setup).

Les \textbf{maintenances} sont organisées en types, chacun défini par :
\begin{itemize}
	\item Une \textbf{périodicité} : intervalle entre deux exécutions (0 pour urgente, $\infty$ pour non applicable)
	\item Un \textbf{ensemble d'opérations} requises
	\item Une \textbf{priorité/pondération} pour la fonction objectif
\end{itemize}

Chaque \textbf{opération} spécifie :
\begin{itemize}
	\item Une \textbf{durée} $dur_o$ (variable selon le type de train)
	\item Des \textbf{relations de précédence} avec d'autres opérations
	\item Des \textbf{exigences en ressources} : type de technicien, équipements requis, compatibilité avec certains types de voies
\end{itemize}

Les \textbf{ressources} se répartissent en trois catégories :
\begin{itemize}
	\item \textbf{Voies} : caractérisées par leur type et étage, ne peuvent accueillir qu'une opération à la fois. Certains équipements ne sont disponibles que sur certaines voies.
	\item \textbf{Techniciens} : organisés par types (spécialités) avec un nombre limité $N_k$ par type et des contraintes de temps de travail ($T_{min,k}$, $T_{max,k}$ par jour).
	\item \textbf{Équipements} : quantité totale et capacité simultanée maximale $C_e$ (nombre d'opérations pouvant l'utiliser simultanément).
\end{itemize}

Les \textbf{paramètres temporels} incluent :
\begin{itemize}
	\item L'horizon de planification (plusieurs jours)
	\item L'unité temporelle (minutes ou périodes)
	\item Deux temps de setup : $\sigma_{op}$ (entre opérations d'un même train) et $\sigma_{train}$ (pour les changements de trains, plus long)
\end{itemize}

\subsection{Variables de décision}

Pour résoudre le problème, on doit déterminer :
\begin{itemize}
	\item $s_o, f_o$ : temps de début et de fin de chaque opération $o$
	\item $s_m, f_m$ : temps de début et de fin de chaque maintenance $m$
	\item $x_o$ : voie assignée à chaque opération $o$
\end{itemize}

Selon la variante du problème, d'autres variables peuvent être nécessaires : allocation spécifique de techniciens aux opérations, activation/désactivation de maintenances optionnelles, horaires de travail détaillés des techniciens.

\subsection{Contraintes}

Une solution valide doit respecter les contraintes suivantes :

\textbf{Contraintes temporelles :}
\begin{itemize}
	\item Cohérence : $f_o = s_o + dur_o$ pour chaque opération $o$
	\item Fenêtres : les maintenances dans les fenêtres de présence des trains
	\item Précédence : $s_o \geq f_{pred(o)} + \sigma_{op}$ si l'opération $o$ a un prédécesseur $pred(o)$
	\item Hiérarchie : opérations dans la fenêtre de leur maintenance
\end{itemize}

\textbf{Contraintes de ressources :}
À tout instant, le nombre de techniciens de type $k$ utilisés $\leq N_k$ et l'utilisation simultanée d'équipement $e$ $\leq C_e$. Les techniciens respectent les temps min/max par jour et ne peuvent effectuer qu'une tâche à la fois (ubiquité).

\textbf{Contraintes de compatibilité :}
Équipements disponibles sur les voies choisies, types de techniciens correspondant aux besoins des opérations.

\textbf{Contraintes de non-chevauchement :}
Deux opérations $o_1$ et $o_2$ sur la même voie (i.e., $x_{o_1} = x_{o_2}$) doivent être disjointes :
\begin{equation}
	f_{o_1} + \sigma \leq s_{o_2} \text{ ou } f_{o_2} + \sigma \leq s_{o_1}
\end{equation}
où $\sigma = \sigma_{op}$ si même train, $\sigma = \sigma_{train}$ sinon.

\subsection{Objectif}

L'objectif est de minimiser le retard pondéré :
\begin{equation}
	\min \sum_{m} w_m \cdot \max(0, f_m - \delta_m)
\end{equation}
où $w_m$ est un poids reflétant la priorité de la maintenance $m$, et $\delta_m$ sa deadline (calculée comme la dernière fois où elle a été effectuée plus sa périodicité). Les maintenances non planifiées peuvent recevoir une forte pénalité selon la variante du problème.
